\documentclass[11pt]{amsart}
\usepackage[margin=1.15in]{geometry}
\usepackage{amsmath,amssymb,amsthm}
\usepackage{booktabs}
\usepackage{graphicx}
\usepackage[hidelinks]{hyperref}
\usepackage{microtype}

\newtheorem{theorem}{Theorem}[section]
\newtheorem{proposition}[theorem]{Proposition}
\newtheorem{lemma}[theorem]{Lemma}
\theoremstyle{remark}
\newtheorem{remark}[theorem]{Remark}
\theoremstyle{definition}
\newtheorem{definition}[theorem]{Definition}

\newcommand{\Q}{\mathbb{Q}}
\newcommand{\R}{\mathbb{R}}
\newcommand{\ran}{r_{\mathrm{an}}}
\DeclareMathOperator{\supp}{supp}

\title[Optimal windowed test functions for analytic rank bounds]{Optimal windowed test functions for bounding\\ analytic ranks of elliptic curves}
\author{Sabrina Hackfort}
\address{Independent researcher (MEGALODON Research), Germany}
\email{meg.depth@proton.me}
\thanks{ORCID: 0009-0007-7458-2487}
\date{October 2026 \quad (draft)}

\begin{document}

\begin{abstract}
Assuming the Generalized Riemann Hypothesis, the explicit formula for the $L$-function of an elliptic curve $E/\Q$ bounds the analytic rank of $E$ in terms of the traces $a_p$ for $p < e^{2L}$, for any admissible test function supported on a window $[-L,L]$. In the literature this test function is fixed in advance, typically to a Fej\'er-type kernel. We instead choose, for each individual curve, the test function that minimises the resulting bound within a finite-dimensional space supported on $[-L,L]$. The minimiser is the solution of a linear system with the Gram matrix of the Weil quadratic form of $E$, and the optimal bound equals the reciprocal of the corresponding reproducing kernel at the central point, which we call the \emph{radix Riemanniana} of $E$. Because the Fej\'er-type kernel lies in this space, the optimised bound is never weaker. We certify every bound with ball arithmetic, check the prime-side evaluation against the zero side computed independently from the zeros of $L(E,s)$, and benchmark against the fixed kernel on $1642$ curves from Cremona's tables, $2500$ curves from the Stein--Watkins database and a set of curves with conductors up to $3\cdot 10^{16}$. Depending on the sample and on the stopping criterion, the number of primes required to pin down the analytic rank decreases by $8$--$23\%$ in aggregate. For individual curves of large conductor the decrease ranges from $0\%$ to $43\%$.
\end{abstract}

\maketitle

\section{Introduction}

Let $E/\Q$ be an elliptic curve of conductor $N$ and let $\ran(E)$ be the order of vanishing of $L(E,s)$ at the central point. Under the Generalized Riemann Hypothesis (GRH), Mestre's explicit-formula method~\cite{Mestre} bounds $\ran(E)$ from above by a quantity computed from the traces of Frobenius $a_p$ for $p$ below a cutoff. Bober~\cite{Bober} made this practical with a Fej\'er-type test function, and Klagsbrun, Sherman and Weigandt~\cite{KSW} used the same function to show that the Elkies curve has rank $28$ subject only to GRH. In these works the shape of the test function is fixed, and only its support (the cutoff) is varied.

This note asks a simple question: \emph{for a given curve and a given cutoff, which test function gives the best bound?} Within any finite-dimensional space of admissible functions supported on $[-L,L]$ the answer is explicit. The explicit formula expresses the bound as a ratio of a quadratic form (the Weil form of $E$, evaluated through the $a_p$) to the square of a linear functional. The minimiser is therefore obtained from one linear solve with a Gram matrix. The optimal value is the reciprocal of the reproducing kernel at $0$ of the space equipped with the Weil inner product.

\subsection*{What is known, and what is new}
The Weil quadratic form restricted to functions supported on a finite interval, and its Gram matrices with respect to a basis, have been studied for the Riemann zeta function by Yoshida~\cite{Yoshida}, by Connes, Consani and Moscovici~\cite{CCM}, by Suzuki~\cite{Suzuki} via screw functions and de Branges spaces, and recently by Shi~\cite{Shi} as finite ``Prime--Weil matrices''. Reproducing-kernel methods for Fourier optimisation problems attached to \emph{families} of $L$-functions, where the input is the limiting one-level density rather than data of an individual $L$-function, are due to Carneiro, Chirre and Milinovich~\cite{CCMil}. To the best of our knowledge, which rests on a literature search that is not exhaustive, the optimisation of the test function \emph{for an individual elliptic curve, using its actual $a_p$}, with certified output and a systematic comparison against the fixed kernel, has not been carried out before. That optimisation is the contribution of this note. The underlying theory is classical. The contribution is the per-curve optimisation, its certification, and the measurements.

\subsection*{Results in brief}
\begin{itemize}
\item The optimised bound is never weaker than the fixed-kernel bound at the same cutoff (Proposition~\ref{prop:dominate}). This is a consequence of the construction, not an empirical finding. The empirical question is \emph{how much} it gains.
\item Every reported bound is certified with ball arithmetic. No certified bound in our tests falls below the known rank.
\item The prime-side value of the Weil form agrees with the zero-side value computed from the zeros of $L(E,s)$. The discrepancy halves with each doubling of the height up to which zeros are included (Section~\ref{sec:zeros}).
\item The number of primes needed to determine the analytic rank (Section~\ref{sec:bench}) decreases in aggregate by $8.3$--$15.2\%$ (Cremona, $N<4.1\cdot 10^5$) and by $17.0$--$22.7\%$ (Stein--Watkins, $4.4\cdot10^6\le N\le 7.5\cdot 10^9$), depending on the stopping criterion. For individual curves of conductor up to $5\cdot 10^{15}$ the decrease ranges from $0\%$ to $43\%$.
\end{itemize}

\section{The explicit formula and rank bounds}

Write $L(E,s)=\sum_{n\ge1} a_n n^{-s-1/2}$ so that the central point is $s=1/2$, and put
\[
\Lambda(E,s) = N^{s/2}(2\pi)^{-s}\,\Gamma(s+\tfrac12)\,L(E,s), \qquad \Lambda(E,s)=\pm\Lambda(E,1-s).
\]
For $p\nmid N$ let $\alpha_p,\beta_p$ be defined by $\alpha_p+\beta_p=a_p/\sqrt p$, $\alpha_p\beta_p=1$, and set $c_{p^k}=\alpha_p^k+\beta_p^k$. For $p\mid N$ set $c_{p^k}=(a_p/\sqrt p)^k$ (so $c_{p^k}=0$ for additive reduction). Then $-\frac{L'}{L}(E,s)=\sum_{n\ge2}\Lambda(n)c_n n^{-s}$ in the normalisation above.

Let $f\colon\R\to\R$ be even, bounded and supported in $[-L,L]$. Put $g=f*f$, so that $g$ is even and supported in $[-2L,2L]$, and
\[
\hat f(r)=\int_{\R} f(x)e^{irx}\,dx,\qquad h(r)=\hat g(r)=\hat f(r)^2 .
\]
Writing the nontrivial zeros as $\rho=\tfrac12+i\gamma$ (with multiplicity), the explicit formula reads
\begin{equation}\label{eq:explicit}
\sum_{\gamma} h(\gamma) \;=\; g(0)\log\frac{N}{4\pi^2}
\;+\;2\Big[-\gamma_{\mathrm E}\,g(0)+\int_0^{\infty}\frac{g(0)-g(u)}{e^{u}-1}\,du\Big]
\;-\;2\sum_{n\ge 2}\frac{\Lambda(n)\,c_n}{\sqrt n}\,g(\log n),
\end{equation}
where $\gamma_{\mathrm E}$ is Euler's constant. Only prime powers $n<e^{2L}$ contribute. We denote the right-hand side by $W_E(f)$.

\begin{lemma}\label{lem:bound}
Assume GRH for $L(E,s)$. Then for every nonzero $f$ as above,
\[
\ran(E)\;\le\;\frac{W_E(f)}{\hat f(0)^2}.
\]
\end{lemma}
\begin{proof}
Under GRH every $\gamma$ is real. Since $f$ is real and even, $\hat f$ is real, so $h(\gamma)=\hat f(\gamma)^2\ge0$. The central zero contributes $\ran(E)\,\hat f(0)^2$ to the left-hand side of~\eqref{eq:explicit}, and every other term is nonnegative.
\end{proof}

\begin{remark}
For $f=\mathbf 1_{[-L,L]}$ one has $g(u)=(2L-|u|)_+$ and $h(r)=4\sin^2(Lr)/r^2$. Up to scaling, this is the Fej\'er-type test function used in~\cite{Bober,KSW}. We call the resulting bound $B_{\mathrm{std}}(L)$.
\end{remark}

\section{The optimal window kernel}

The map $f\mapsto W_E(f)$ is a quadratic form, $W_E(f)=\mathcal B_E(f,f)$, where $\mathcal B_E$ is obtained from~\eqref{eq:explicit} by replacing $g=f*f$ with the symmetrised correlation of two functions. Fix the basis
\[
b_k(x)=P_{2k}(x/L)\,\mathbf 1_{[-L,L]}(x),\qquad 0\le k<K,
\]
of even polynomials on the window, where $P_m$ is the Legendre polynomial of degree $m$. Let $G_{ij}=\mathcal B_E(b_i,b_j)$ be the Gram matrix and $v_i=\hat b_i(0)=\int b_i$, so that $v=(2L,0,\dots,0)$. For $f=\sum c_k b_k$ we have $W_E(f)=c^{\mathsf T}Gc$ and $\hat f(0)=v^{\mathsf T}c$.

\begin{proposition}\label{prop:opt}
If $G$ is positive definite, the minimum of $W_E(f)/\hat f(0)^2$ over nonzero $f$ in $\mathrm{span}(b_0,\dots,b_{K-1})$ equals
\[
B_{\mathrm{rad}}(L)=\frac{1}{v^{\mathsf T}G^{-1}v},
\]
and is attained at $c^\ast=G^{-1}v/(v^{\mathsf T}G^{-1}v)$.
\end{proposition}
\begin{proof}
Minimise $c^{\mathsf T}Gc$ subject to $v^{\mathsf T}c=1$ (Lagrange multipliers, or Cauchy--Schwarz in the inner product $\langle x,y\rangle=x^{\mathsf T}Gy$).
\end{proof}

Since $v=\Phi(0)$, the quantity $v^{\mathsf T}G^{-1}v$ is the value at the central point of the following kernel.

\begin{definition}[Radix Riemanniana]\label{def:radix}
Let $\Phi(z)=(\hat b_0(z),\dots,\hat b_{K-1}(z))^{\mathsf T}$. For an elliptic curve $E/\Q$ and a window $[-L,L]$, the \emph{radix Riemanniana} of $E$ (of order $K$) is the kernel
\[
\mathcal R_{E,L}(w,z)\;=\;\Phi(w)^{\ast}\,G^{-1}\,\Phi(z),
\]
that is, the reproducing kernel of the space $\{\hat f: f\in\mathrm{span}(b_k)\}$ equipped with the Weil inner product $\mathcal B_E$. It is built from the Gamma factor, the conductor and the primes $p<e^{2L}$ alone, without reference to the zeros of $L(E,s)$. The optimal bound of Proposition~\ref{prop:opt} is
\[
B_{\mathrm{rad}}(L)=\frac{1}{\mathcal R_{E,L}(0,0)}.
\]
\end{definition}

\begin{remark}
As $K\to\infty$ the radix Riemanniana is the window analogue, for $L(E,s)$, of the reproducing kernels of the Hilbert spaces of entire functions attached to the Weil form of $\zeta$ in~\cite{Suzuki}. The name and the construction from prime data go back to the author's work on Weil positivity for $\zeta$.
\end{remark}

\begin{proposition}\label{prop:dominate}
For every $L$ with $G$ positive definite, $B_{\mathrm{rad}}(L)\le B_{\mathrm{std}}(L)$.
\end{proposition}
\begin{proof}
$\mathbf 1_{[-L,L]}=b_0$ lies in the span, so $B_{\mathrm{std}}(L)$ is one of the values minimised over in Proposition~\ref{prop:opt}.
\end{proof}

Positive definiteness of $G$ holds under GRH except in degenerate cases, since $W_E(f)=\sum_\gamma\hat f(\gamma)^2$. In practice we simply attempt a Cholesky factorisation and discard the window if it fails. Note that \emph{the certificate in Section~\ref{sec:cert} does not depend on optimality or on positive definiteness}: it evaluates $W_E(f)$ for the specific $f$ that was found.

\section{Certification}\label{sec:cert}

The optimisation is carried out in double precision. The resulting coefficients are converted to exact rationals $c_k\in\Q$, and $L$ is taken to be an exact decimal. For this fixed $f=\sum c_k b_k$ we evaluate $W_E(f)/\hat f(0)^2$ rigorously in ball arithmetic (Arb~\cite{Arb} via \texttt{python-flint}, $256$-bit precision):
\begin{enumerate}
\item On $[0,2L]$, $g=f*f$ is a polynomial with rational coefficients. We compute it exactly in $\Q[u]$ and obtain $\hat f(0)$ exactly.
\item The archimedean integral is split at $u=1/4$. On $[0,1/4]$ we integrate $(g(0)-g(u))/u$ against the Bernoulli expansion $u/(e^u-1)=\sum_n B_n u^n/n!$, truncated after $80$ terms. The tail is bounded with $|B_n|/n!\le 4(2\pi)^{-n}$. On $[1/4,2L]$ the integrand is analytic and is integrated rigorously by Arb's certified Gauss--Legendre integration. The constants $\gamma_{\mathrm E}$, $\pi$ and the logarithms are balls.
\item The prime sum runs over all prime powers $n<e^{2L}$. The quantities $\log n$, $\sqrt n$ and $c_n$ are balls computed from the exact integers $a_p$. If the ball for $\log n$ contains $2L$, the window is rejected.
\end{enumerate}
The output is a ball containing $W_E(f)/\hat f(0)^2$. A statement $\ran(E)<r+1$ is recorded only if the upper end of the ball is below $r+1$. The $a_p$ are computed with PARI/GP~\cite{PARI}.

\section{Validation against the zeros}\label{sec:zeros}

Identity~\eqref{eq:explicit} lets us check the whole pipeline (Rust optimiser $\to$ rational coefficients $\to$ certified prime side) against an independent computation. We compute the zeros $\gamma$ of $L(E,s)$ up to height $T$ with PARI's \texttt{lfunzeros} and evaluate the zero side $\ran\hat f(0)^2+\sum_{0<|\gamma|\le T}\hat f(\gamma)^2$ for the same certified $f$. Since every omitted term is nonnegative, the truncated zero side must lie below the prime side, and the difference must tend to $0$ as $T\to\infty$.

\begin{table}[h]
\centering\small
\begin{tabular}{lcccccc}
\toprule
curve & $\ran$ & $L$ & prime side (certified) & zero side, $|\gamma|\le T$ & $T$ & difference\\
\midrule
37a1     & 1 & 0.40 & 1.9967148280 & 1.9923507969 & 600 & $4.4\cdot10^{-3}$\\
389a1    & 2 & 0.65 & 2.9926821162 & 2.9903944027 & 600 & $2.3\cdot10^{-3}$\\
5077a1   & 3 & 0.90 & 3.9514652620 & 3.9497072474 & 600 & $1.8\cdot10^{-3}$\\
234446a1 & 4 & 1.15 & 4.9323755046 & 4.9292024883 & 300 & $3.2\cdot10^{-3}$\\
\bottomrule
\end{tabular}
\caption{Certified prime side versus zero side, both normalised by $\hat f(0)^2$, for the optimised test function at the cutoff where the bound first drops below $\ran+1$.}
\label{tab:zeros}
\end{table}

For 37a1 the difference is $8.1\cdot10^{-3}$, $4.4\cdot10^{-3}$ and $2.4\cdot10^{-3}$ at $T=300$, $600$ and $1200$ ($446$, $1024$ and $2313$ zeros). It roughly halves with each doubling of $T$, consistent with a truncation tail of order $(\log T)/T$.

\section{Benchmarks}\label{sec:bench}

For each curve we scan $L$ on a grid of step $0.05$ with $K=10$ basis functions. For both $B_{\mathrm{std}}$ and $B_{\mathrm{rad}}$ we record the first $L$ at which the certified bound drops below $\ran+1$ (criterion~A), and the first $L$ at which it drops below $\ran+2$ (criterion~B; together with the root number, this already determines the parity of $\ran$). Cost is measured as the number of primes $p<e^{2L}$ used. Because of the grid, prime counts at the threshold are quantised. Ties are therefore frequent and slightly understate the gain.

\subsection{Cremona's tables}
We take one curve per isogeny class, sampled up to $150$ per rank from three blocks of Cremona's tables~\cite{Cremona} ($61\le N\le 409989$, $1642$ curves, ranks $0$--$3$), with the cutoff $L\le 4.5$ ($1028$ primes). Table~\ref{tab:cremona} reports the results. In total $6449$ bounds were certified, and none lies below the known rank.

\begin{table}[h]
\centering\small
\begin{tabular}{lccccc}
\toprule
 & curves decided (std / rad) & rad fewer / equal / more & primes saved\\
\midrule
criterion A ($<\ran+1$) & 1573 / 1592 & 578 / 1064 / 0 & 15.2\%\\
criterion B ($<\ran+2$) & 1642 / 1642 & 263 / 1379 / 0 & \phantom{1}8.3\%\\
\bottomrule
\end{tabular}

\smallskip
\begin{tabular}{lcccc}
\toprule
criterion A, by rank & 0 & 1 & 2 & 3\\
\midrule
curves (both decided) & 381 & 450 & 450 & 292\\
primes saved & 18.8\% & 4.3\% & 3.9\% & 4.5\%\\
\bottomrule
\end{tabular}
\caption{Cremona sample. ``Primes saved'' is $1-\sum P_{\mathrm{rad}}/\sum P_{\mathrm{std}}$ over curves decided by both methods.}
\label{tab:cremona}
\end{table}

\subsection{The Stein--Watkins database}
From the Stein--Watkins database~\cite{SW} (files of the 2011 distribution) we sample $250$ curves per (analytic) rank $0,1,2,3,\ge4$ from two groups: general conductors $4.4\cdot10^6\le N\le9.7\cdot10^7$ and prime conductors $3.6\cdot10^9\le N\le7.5\cdot10^9$, giving $2500$ curves in total. We check every conductor and root number with PARI. The ranks in the database are Watkins' numerically determined analytic ranks. The cutoff is $L\le5.75$ ($9592$ primes). In total $9105$ bounds were certified, and none contradicts the database rank.

\begin{table}[h]
\centering\small
\begin{tabular}{lccc}
\toprule
 & decided (std / rad) & rad fewer / equal / more & primes saved\\
\midrule
criterion A & 2059 / 2096 & 1306 / 1194 / 0 & 17.0\%\\
criterion B & 2466 / 2484 & \phantom{1}889 / 1611 / 0 & 22.7\%\\
\bottomrule
\end{tabular}

\smallskip
\begin{tabular}{lccccc}
\toprule
criterion B, by rank & 0 & 1 & 2 & 3 & $\ge4$\\
\midrule
curves (both decided) & 466 & 500 & 500 & 500 & 500\\
primes saved & 24.7\% & 15.1\% & 5.7\% & 3.1\% & 2.7\%\\
\bottomrule
\end{tabular}
\caption{Stein--Watkins sample.}
\label{tab:sw}
\end{table}

\subsection{Large conductors and high rank}
Table~\ref{tab:big} lists every curve of an additional test set, including the cases without gain. Prime counts up to $e^{16}$ were handled by streaming the prime sum.

\begin{table}[h]
\centering\small
\begin{tabular}{lrcccc}
\toprule
curve & $N$ & rank & criterion & primes std / rad & saved\\
\midrule
twist of 11a1, $d=10009$     & $1101980891$         & 1 & B & 329 / 302 & 8\%\\
twist of 389a1, $d=10009$    & $38970051509$        & 1 & B & 3523 / 3523 & 0\%\\
twist of 11a1, $d=100049$    & $110107826411$       & 0 & B & 2466 / 2065 & 16\%\\
twist of 5077a1, $d=10009$   & $508614271237$       & 1 & B & 1220 / 1019 & 16\%\\
twist of 389a1, $d=100049$   & $3893813133989$      & 0$^{\ast}$ & B & 17885 / 13613 & 24\%\\
twist of 5077a1, $d=100049$  & $50819766789877$     & 1 & B & 1220 / 1220 & 0\%\\
twist of 389a1, $d=1000033$  & $389025674423621$    & 0$^{\ast}$ & B & 102084 / 77395 & 24\%\\
twist of 5077a1, $d=1000033$ & $5077335087528853$   & 1 & B & 340436 / 195040 & 43\%\\
twist of 11a1, $d=1000033$   & $11000726011979$     & 0 & B & \multicolumn{2}{c}{undecided by both for $L\le7.5$}\\
\midrule
5187563742a1                 & $5187563742$         & 6 & A & 11 / 10 & 9\%\\
$[1,-1,0,-106384,13075804]$  & $249649566346838$    & 8 & A & 21 / 18 & 14\%\\
$[1,-1,0,-135004,97151644]$  & $32107342006814614$  & 9 & A & 32 / 29 & 9\%\\
\bottomrule
\end{tabular}
\caption{Additional curves. Ranks are lower bounds from PARI's \texttt{ellrank} (points). $^{\ast}$Here the 2-descent leaves rank $0$ or $2$ open. The root number $+1$ together with the certified bound $\ran<2$ gives $\ran=0$ under GRH, and hence rank $0$ by Kolyvagin~\cite{Kolyvagin}.}
\label{tab:big}
\end{table}

\section{Limitations}
\begin{itemize}
\item All bounds assume GRH for $L(E,s)$ and concern the analytic rank. Statements about the algebraic rank additionally use points, the root number and, for $\ran\le1$, the theorems of Gross--Zagier~\cite{GZ} and Kolyvagin~\cite{Kolyvagin}.
\item The basis size is fixed at $K=10$ and the cutoff grid at step $0.05$. Larger $K$ can only improve $B_{\mathrm{rad}}$, and a finer grid would turn some ties into gains. We have not optimised either choice.
\item The certification is currently performed in Python on top of Arb. The optimiser is written in Rust.
\item Our samples cover conductors up to $\sim3\cdot10^{16}$ and ranks up to $9$. We have not tested curves of record rank such as the Elkies curve, whose $a_p$ are costly to compute. That is where savings in the number of primes would matter most.
\end{itemize}

\section*{Data and code availability}
Code (Rust optimiser, certification scripts), the sampled curve lists and all certificates are available at \url{https://github.com/sectio-aurea-q/radix-rank} (archived on Zenodo, DOI: \href{https://doi.org/10.5281/zenodo.23299461}{10.5281/zenodo.23299461}). Running \texttt{./reproduce.sh quick} re-derives and re-certifies the bounds for seven reference curves in about one minute.

\section*{Acknowledgements}
An AI assistant (Claude, Anthropic) was used in preparing this work. All results were verified computationally as described above.

\begin{thebibliography}{99}
\bibitem{Arb} F.~Johansson, \emph{Arb: efficient arbitrary-precision midpoint-radius interval arithmetic}, IEEE Trans.\ Comput.\ \textbf{66} (2017), 1281--1292.
\bibitem{Bober} J.~W.~Bober, \emph{Conditionally bounding analytic ranks of elliptic curves}, ANTS X, Open Book Ser.\ \textbf{1} (2013), 135--144; arXiv:1112.1503.
\bibitem{CCMil} E.~Carneiro, A.~Chirre, M.~B.~Milinovich, \emph{Hilbert spaces and low-lying zeros of $L$-functions}, arXiv:2109.10844.
\bibitem{CCM} A.~Connes, C.~Consani, H.~Moscovici, \emph{Zeta zeros and prolate wave operators}, arXiv:2310.18423.
\bibitem{Cremona} J.~E.~Cremona, \emph{Algorithms for Modular Elliptic Curves}, 2nd ed., Cambridge Univ.\ Press, 1997; tables at \url{https://github.com/JohnCremona/ecdata}.
\bibitem{GZ} B.~H.~Gross, D.~B.~Zagier, \emph{Heegner points and derivatives of $L$-series}, Invent.\ Math.\ \textbf{84} (1986), 225--320.
\bibitem{KSW} Z.~Klagsbrun, T.~Sherman, J.~Weigandt, \emph{The Elkies curve has rank 28 subject only to GRH}, Math.\ Comp.\ \textbf{88} (2019), 837--846; arXiv:1606.07178.
\bibitem{Kolyvagin} V.~A.~Kolyvagin, \emph{Euler systems}, in: The Grothendieck Festschrift, Vol.~II, Birkh\"auser, 1990, 435--483.
\bibitem{Mestre} J.-F.~Mestre, \emph{Formules explicites et minorations de conducteurs de vari\'et\'es alg\'ebriques}, Compositio Math.\ \textbf{58} (1986), 209--232.
\bibitem{PARI} The PARI~Group, \emph{PARI/GP}, Univ.\ Bordeaux, \url{https://pari.math.u-bordeaux.fr}.
\bibitem{Shi} Y.~Shi, \emph{Construction of finite Hilbert--P\'olya matrices from Weil's explicit formula}, arXiv:2609.04908.
\bibitem{SW} W.~A.~Stein, M.~Watkins, \emph{A database of elliptic curves---first report}, ANTS V, Lecture Notes in Comput.\ Sci.\ \textbf{2369} (2002), 267--275.
\bibitem{Suzuki} M.~Suzuki, \emph{Weil's quadratic form via the screw function}, arXiv:2606.09096.
\bibitem{Yoshida} H.~Yoshida, \emph{On Hermitian forms attached to zeta functions}, in: Zeta Functions in Geometry, Adv.\ Stud.\ Pure Math.\ \textbf{21} (1992), 281--325.
\end{thebibliography}

\end{document}
